\documentclass{article}

\usepackage[utf8]{inputenc}
\usepackage{amsmath,amssymb,amsfonts}
\usepackage{graphicx}
\usepackage{hyperref}
\usepackage{booktabs}

\input{macros}

\title{
Hyperbolic Spectral Projectors and the Quaternionic Polar Chart on Split Octonions:
A Formally Verified Framework in Lean 4
}

\author{Nikolay Goutev \and Dimitar Tonev}

\begin{document}

\maketitle

\begin{abstract}
We formalize a rigorous hyperbolic spectral framework on the split octonions in Lean 4.
By utilizing a quaternionic-hyperbolic polar chart on the invertible sectors, we define the radial coordinate $\rho$ and rapidity $\eta$, and extract a canonical pure-imaginary generator $J_g$. We formally verify the global polar decomposition $X = \rho(u, 0) e^{\eta J_g}$ using exact algebraic identities, completely avoiding $C^*$-algebra approximations. We further extend this to construct twin-wave spectral projectors $p_\pm = \frac{1}{2}(1 \pm J_g)$ that capture the reciprocal spectral sheets, and we establish the bridging relations between the hyperbolic and elliptic structures.

Our framework strictly adheres to native mathlib structures, employing pure algebraic proofs and discarding any legacy witness-based axioms. This solidifies the foundational geometry required for the algebraic treatment of topological quantum field theories and operator algebras.
\end{abstract}

\section*{Claim Hierarchy and Verification Status}
To maintain a strict separation between formally verified facts and mathematical hypotheses, every core statement in this manuscript is assigned to one of four verification classes. The classification is summarized in Table~\ref{tab:claim_hierarchy}.

\begin{table}[h]
\centering
\begin{tabular}{llp{7cm}}
\toprule
\textbf{Class} & \textbf{Meaning} & \textbf{Current Status in Repository} \\
\midrule
\textbf{Theorem} & Compiled Lean theorem with published \texttt{\#print axioms} output. & Quaternionic polar chart \texttt{polar\_decomposition} (verified), $J_g^2 = 1$ property (verified), Twin-wave projectors (pending). \\
\textbf{Implementation invariant} & Enforced and regression-tested. & Pure mathlib import paths, absence of \texttt{sorry} and witness gaps. \\
\textbf{Benchmark claim} & Supported by reproducible log metrics. & Lean 4 kernel compilation times and exact term sizes. \\
\textbf{Mathematical hypothesis} & Standard unverified assumption. & Extensions to infinite-dimensional Tomita modular conjugations via $L^2(\mathrm{Cl}^+(1,3))$. \\
\bottomrule
\end{tabular}
\caption{Classification of manuscript claims by verification standard.}
\label{tab:claim_hierarchy}
\end{table}
% Section 01_introduction

\section{Calibrated CMOS Quasi-Deviance}
\label{sec:cmos_quasideviance}

To develop a principled anomaly detector for RAW CMOS sequences, we begin with the underlying physical variance law of the sensor. The total heteroscedastic variance of a pixel, under standard operation, is modeled by the polynomial:
\begin{equation}
V(\mu) = c_0 + c_1\mu + c_2\mu^2,
\end{equation}
subject to the constraints
\begin{equation}
c_0 > 0, \qquad c_1, c_2 \ge 0.
\end{equation}
Here, $c_0$ captures the constant read noise floor, $c_1\mu$ models the Poisson-distributed photon shot noise, and $c_2\mu^2$ accounts for multiplicative gain variations or flat-fielding uncertainties at high signal levels.

Given this variance function, we construct the quasi-deviance observation model as the foundational physical metric for the estimator. The quasi-deviance $D_V(x | \mu)$ between an observation $x$ and an expected background $\mu$ is defined via the integral:
\begin{equation}
D_V(x|\mu) = 2\int_\mu^x \frac{x-u}{V(u)}\,du.
\end{equation}

By construction, this quasi-deviance locally scales the squared prediction error by the physical variance, smoothly interpolating between Gaussian geometry in the read-noise limit, Poisson geometry in the shot-noise limit, and Itakura--Saito (Gamma) geometry in the high-signal limit. This provides a robust, physics-informed distance measure that governs the admission responsibility of incoming pixels without requiring ad hoc variance stabilization transforms.

% Section 03_mixed_bregman_geometry

% Section 04_recurrent_estimator

% Section 05_formal_verification

% Section 06_implementation

% Section 07_experiments

% Section 08_limitations

% Section 09_conclusion


\bibliographystyle{plain}
\bibliography{references}

\appendix
\section{Artifacts}
\input{artifacts/theorem_status}
\input{artifacts/benchmark_environment}
\input{artifacts/reproducibility_manifest}

\end{document}
